\documentclass[11pt]{article}
\usepackage[margin=1in]{geometry}
\usepackage{amsmath,amssymb,amsthm,mathtools}
\usepackage[T1]{fontenc}
\usepackage{enumitem}
\usepackage{xcolor}
\usepackage{hyperref}

\hypersetup{colorlinks=true, linkcolor=blue!60!black, citecolor=blue!60!black, urlcolor=blue!60!black}

\newtheorem{theorem}{Theorem}[section]
\newtheorem{lemma}[theorem]{Lemma}
\newtheorem{proposition}[theorem]{Proposition}
\newtheorem{corollary}[theorem]{Corollary}
\newtheorem{definition}[theorem]{Definition}
\newtheorem{remark}[theorem]{Remark}

\newcommand{\R}{\mathbb R}
\newcommand{\N}{\mathbb N}
\newcommand{\Poly}{\mathcal P}
\newcommand{\Lin}{\mathcal E}
\newcommand{\Quad}{\mathcal Q}
\newcommand{\Quart}{\mathcal T}
\newcommand{\Span}{\operatorname{span}}
\newcommand{\rank}{\operatorname{rank}}
\newcommand{\codim}{\operatorname{codim}}
\newcommand{\Ann}{\operatorname{Ann}}
\newcommand{\Sym}{\operatorname{Sym}}
\newcommand{\im}{\operatorname{im}}
\newcommand{\kerop}{\operatorname{ker}}
\newcommand{\sos}{\operatorname{sos}}
\newcommand{\td}{\operatorname{totalDegree}}
\newcommand{\defeq}{\mathrel{:=}}
\newcommand{\lean}[1]{\texttt{\detokenize{#1}}}

\title{Blueprint for \lean{QuaternaryQuarticRankSevenNoSpuriousSOCP}}
\author{}
\date{}

\begin{document}
\maketitle

\section{Purpose and exact target}

This blueprint is written for the immutable Lean target in \lean{QuaternaryQuartic.lean}:
\[
\begin{aligned}
\lean{QuaternaryQuarticRankSevenNoSpuriousSOCP} :\quad
\forall (B : \lean{DotForm})\,(p : \lean{Poly})\,(u : \lean{RankSevenVec}),\;&
\lean{IsPositiveDefinite}\;B \to
\lean{IsSOSQuartic}\;p \to{}\\
&\lean{IsAdmissiblePoint}\;u \to
\lean{IsSOCP}\;B\;p\;u \to
\lean{residual}\;p\;u = 0.
\end{aligned}
\]
Here
\[
\Poly=\R[X_1,X_2,X_3],\qquad
\Quad=\{q\in\Poly: \td(q)\le 2\},
\]
so dehomogenized quaternary quadratics form a 10-dimensional real vector space.  The rank-seven factor variable is a map
\[
 u:\{0,\dots,6\}\to \Poly,
\]
with each coordinate in \(\Quad\).  The quadratic map and its linearization are
\[
\sigma(u)=\sum_{i=0}^6 u_i^2,
\qquad
A_u(v)=\sum_{i=0}^6 u_i v_i.
\]
For a fixed quartic target \(p\), the residual is
\[
R=\sigma(u)-p.
\]
The Lean file defines the second-order expression as
\[
H_u(v)=B(\sigma(v),R)+2B(A_u(v),A_u(v)).
\]
Thus \lean{IsSOCP} means
\[
\forall v\in \Quad^7,\quad B(A_u(v),R)=0,
\qquad
\forall v\in \Quad^7,\quad 0\le H_u(v).
\]
The proof below uses only these definitions.  In particular, it does \emph{not} require symmetry of \(B\).  Positive-definiteness is used only in the form
\[
R\ne 0\implies B(R,R)>0.
\]

\begin{remark}[Relation to the homogeneous statement]
The dehomogenized space \(\Quad\) is naturally isomorphic to the homogeneous quadratic space in four variables.  Let
\[
\Lin=\{\ell\in\Poly: \td(\ell)\le 1\}.
\]
Then \(\dim\Lin=4\), \(\dim\Quad=10\), and products \(\ell m\), \(\ell,m\in\Lin\), span \(\Quad\).  This is the affine chart of the homogeneous quaternary-quartic problem.
\end{remark}

\section{Finite-dimensional spaces and notation}

For formalization it is useful to introduce submodules or subtypes:
\[
\Lin=\{a: \td(a)\le 1\},\qquad
\Quad=\{q: \td(q)\le 2\},\qquad
\Quart=\{f: \td(f)\le 4\}.
\]
The required dimensions are
\[
\dim_\R \Lin=4,
\qquad
\dim_\R \Quad=10,
\qquad
\dim_\R \Quart=35.
\]
The proofs only need \(\dim\Lin=4\) and \(\dim\Quad=10\), plus the closure facts
\[
\Lin\cdot\Lin\subseteq\Quad,
\qquad
\Quad\cdot\Quad\subseteq\Quart.
\]
If \(U,W\subseteq\Lin\), write
\[
UW\defeq \Span_\R\{uw:u\in U,w\in W\}\subseteq\Quad,
\]
and write \(\Sym^2 U\) for \(UU\).  If \(E=U\oplus W\), then
\[
\Sym^2 E=(\Sym^2 U)\oplus (UW)\oplus(\Sym^2 W),
\]
with the expected dimensions.  The only cases used later are:
\[
\dim \Sym^2(\R^2)=3,
\qquad
\dim \Sym^2(\R^3)=6,
\qquad
\dim(xA)=\dim A
\]
for a nonzero linear form \(x\notin A\).  These are elementary basis computations.

\section{The residual functional}

Fix \(B,p,u\) satisfying the hypotheses except for the desired conclusion.  Let
\[
R=\sigma(u)-p.
\]
Assume for contradiction that \(R\ne0\).  Define the linear functional
\[
\lambda:\Poly\to\R,
\qquad
\lambda(f)=B(f,R).
\]
The entire proof is driven by the restriction of \(\lambda\) to quartics.

\begin{lemma}[FOCP makes \(\lambda(p)\) negative]\label{lem:lambda-p-negative}
If \(R\ne0\), \(u\) is admissible, \(B\) is positive-definite, and \(u\) is first-order critical, then
\[
\lambda(p)<0.
\]
\end{lemma}

\begin{proof}
Because \(u\) is admissible, the direction \(v=u\) is admissible.  First-order criticality gives
\[
0=B(A_u(u),R)=B(\sigma(u),R)=\lambda(\sigma(u)).
\]
Since \(R=\sigma(u)-p\), we have \(p=\sigma(u)-R\).  Therefore
\[
\lambda(p)=\lambda(\sigma(u))-\lambda(R)=0-B(R,R).
\]
Positive-definiteness and \(R\ne0\) give \(B(R,R)>0\), hence \(\lambda(p)<0\).
\end{proof}

\begin{lemma}[A negative square from the SOS certificate]\label{lem:negative-sos-summand}
Assume \(p\) is SOS in the sense of \lean{IsSOSQuartic}, and \(\lambda(p)<0\).  Then there exists a quadratic polynomial \(q\in\Quad\) such that
\[
\lambda(q^2)<0.
\]
\end{lemma}

\begin{proof}
Write the SOS certificate as
\[
p=\sum_{j\in J} q_j^2,
\qquad q_j\in\Quad.
\]
By linearity,
\[
\lambda(p)=\sum_{j\in J}\lambda(q_j^2).
\]
If every term were nonnegative, the sum would be nonnegative.  Since \(\lambda(p)<0\), at least one term is negative.
\end{proof}

\section{The catalecticant kernel attached to \(\lambda\)}

Define the degree-two catalecticant bilinear form
\[
C_\lambda:\Quad\times\Quad\to\R,
\qquad
C_\lambda(a,b)=\lambda(ab).
\]
Because multiplication is commutative, \(C_\lambda(a,b)=C_\lambda(b,a)\).  Its left radical is
\[
K=K_\lambda
\defeq
\{a\in\Quad: \forall b\in\Quad,\ \lambda(ab)=0\}.
\]
Equivalently, \(K=\ker(\Quad\to\Quad^*)\), \(a\mapsto(b\mapsto\lambda(ab))\).  Put
\[
k=\rank C_\lambda=10-\dim K.
\]

Let
\[
L=L_u\defeq \Span_\R\{u_0,\dots,u_6\}\subseteq\Quad.
\]

\begin{lemma}[FOCP implies \(L\subseteq K\)]\label{lem:L-subset-K}
If \(u\) is first-order critical, then every \(a\in L\) lies in \(K\).  Equivalently,
\[
L\subseteq K.
\]
\end{lemma}

\begin{proof}
Fix an index \(i\) and any quadratic \(q\in\Quad\).  Let \(v\) be the rank-seven direction with \(v_i=q\) and all other coordinates zero.  Then
\[
A_u(v)=u_iq.
\]
First-order criticality gives \(\lambda(u_iq)=0\).  By linearity in the coordinate \(u_i\), every element of their span has the same property.
\end{proof}

\section{Rank-deficient factor points}

The rank-deficient case is easier than in the SDP proof; no orthogonal rotation is needed.

\begin{lemma}[Rank-deficient points have negative curvature unless \(R=0\)]\label{lem:rank-deficient}
Assume \(R\ne0\), \(p\) is SOS, \(u\) is admissible and first-order critical.  If the seven quadratics \(u_i\) are linearly dependent, then \(u\) is not second-order critical.
\end{lemma}

\begin{proof}
By Lemmas~\ref{lem:lambda-p-negative} and~\ref{lem:negative-sos-summand}, choose \(q\in\Quad\) with \(\lambda(q^2)<0\).  Since the \(u_i\) are linearly dependent, choose coefficients \(c_i\), not all zero, with
\[
\sum_i c_i u_i=0.
\]
Define an admissible direction
\[
v_i=c_iq.
\]
Then
\[
A_u(v)=\sum_i u_i(c_iq)=q\sum_i c_i u_i=0.
\]
Also
\[
\sigma(v)=\sum_i(c_iq)^2=\left(\sum_i c_i^2\right)q^2.
\]
The scalar \(\sum_i c_i^2
>0\), so
\[
H_u(v)
=
\lambda(\sigma(v))+2B(0,0)
=
\left(\sum_i c_i^2\right)\lambda(q^2)<0.
\]
This contradicts second-order criticality.
\end{proof}

Thus, under the contradiction assumption \(R\ne0\), a second-order critical point must have
\[
\dim L=7.
\]
Combined with \(L\subseteq K\), this implies
\[
\dim K\ge7,
\qquad
k=10-\dim K\le3.
\]

\section{The syzygy-to-negative-curvature mechanism}

The next lemma is the main reusable Hessian argument.  It converts an algebraic product identity into a negative second derivative direction.

\begin{lemma}[Syzygy gives negative curvature]\label{lem:syzygy-negative-curvature}
Assume \(u\) is admissible and first-order critical.  Let \(L\subseteq K\) be as above.  Suppose there exist a finite index set \(J\), elements
\[
m_j\in L,
\qquad
n_j\in K,
\qquad
s\in L\setminus\{0\},
\qquad
q\in\Quad,
\]
with
\[
\sum_{j\in J}m_jn_j=sq,
\qquad
\lambda(q^2)<0.
\]
If the coordinates \(u_i\) are linearly independent, then \(u\) is not second-order critical.
\end{lemma}

\begin{proof}
Write
\[
m_j=\sum_i \alpha_{j,i}u_i,
\qquad
s=\sum_i\beta_i u_i.
\]
Because \(s\ne0\) and the \(u_i\) are linearly independent, the coefficient vector \(\beta=(\beta_i)_i\) is nonzero, hence
\[
\sum_i\beta_i^2>0.
\]
Define the admissible direction
\[
v_i=\sum_{j\in J}\alpha_{j,i}n_j-\beta_iq.
\]
Then
\[
A_u(v)
=
\sum_i u_iv_i
=
\sum_{j\in J}\left(\sum_i\alpha_{j,i}u_i\right)n_j
-
\left(\sum_i\beta_i u_i\right)q
=
\sum_{j\in J}m_jn_j-sq
=0.
\]
Now compute \(\lambda(\sigma(v))\).  For every \(n_j\in K\) and every quadratic \(r\), one has \(\lambda(n_jr)=0\).  Therefore all terms involving at least one \(n_j\) vanish under \(\lambda\).  Only the \(q^2\)-terms remain:
\[
\lambda(\sigma(v))
=
\sum_i\lambda(v_i^2)
=
\sum_i\beta_i^2\lambda(q^2)<0.
\]
Consequently
\[
H_u(v)=\lambda(\sigma(v))+2B(A_u(v),A_u(v))
=\left(\sum_i\beta_i^2\right)\lambda(q^2)<0,
\]
contradicting second-order criticality.
\end{proof}

\section{Low-rank catalecticants}

This is the only genuinely structural input.  It is best isolated as a separate formalization block.

\subsection{Apolar support reduction}

Let \(E\) be the four-dimensional affine-linear space \(\Lin\).  The quadratic space \(\Quad\) is \(\Sym^2E\) after homogenization.  Given \(\lambda\) on quartics, define
\[
\Ann_1(\lambda)=\{a\in E: \forall q\in\Sym^3E,\ \lambda(aq)=0\}.
\]
Equivalently, since cubics are spanned by products of a linear form and a quadratic,
\[
a\in\Ann_1(\lambda)
\quad\Longleftrightarrow\quad
\forall e\in E,
\ ae\in K.
\]

\begin{lemma}[Low catalecticant rank forces few essential variables]\label{lem:low-rank-support}
Let \(C_\lambda:\Sym^2E\to(\Sym^2E)^*\) have rank \(k\le3\).  Then there is a direct sum decomposition
\[
E=W\oplus A
\]
such that
\[
\dim W\le k,
\qquad
\dim A\ge4-k,
\qquad
A E\subseteq K.
\]
Equivalently, \(\lambda\) only sees the \(W\)-part: if a quadratic has at least one factor in \(A\), then it lies in \(K\), and any square-pairing involving it is zero under \(\lambda\).
\end{lemma}

\begin{proof}[Proof idea for formalization]
Homogenize and use the apolar Artinian Gorenstein algebra
\[
A_\lambda=S/\Ann(\lambda),
\qquad S=\R[E^*].
\]
Its Hilbert function has the symmetric form
\[
h(A_\lambda)=(1,h_1,h_2,h_1,1),
\qquad
h_2=\rank C_\lambda=k.
\]
The number \(h_1
\) is the number of essential variables, namely
\[
h_1=4-\dim\Ann_1(\lambda).
\]
It remains to show \(h_1\le k\) for \(k\le3\).

For \(k\le2\), Macaulay growth in degree two gives
\[
h_3\le h_2^{\langle2\rangle}\le2.
\]
By Gorenstein symmetry, \(h_1=h_3\le2\).

For \(k=3\), suppose for contradiction that \(h_1=4\).  Then
\[
h=(1,4,3,4,1).
\]
Choose a nonzero linear element \(\ell\in(A_\lambda)_1\).  The standard exact sequence
\[
0\to A_\lambda/(0:\ell)(-1)
\xrightarrow{\cdot\ell}
A_\lambda
\to
A_\lambda/\ell A_\lambda
\to0
\]
has first term Gorenstein of socle degree three, hence Hilbert function
\[
(1,b,b,1)
\]
for some \(1\le b\le3\).  Therefore the quotient \(A_\lambda/\ell A_\lambda\), a standard graded quotient in at most three variables, has Hilbert function
\[
(1,3,3-b,4-b,0).
\]
But this is not an \(O\)-sequence: Macaulay's degree-two bound says
\[
4-b\le(3-b)^{\langle2\rangle},
\]
and for \(b=1,2,3\) the right side is respectively \(2,1,0\), always smaller than \(4-b\).  Contradiction.  Thus \(h_1\le3\).

Choose \(A=\Ann_1(\lambda)\subseteq E\) and any complement \(W\).  Since \(\dim W=h_1\le k\), and since every element of \(A\) annihilates all degree-three products, one gets \(AE\subseteq K\).
\end{proof}

\begin{remark}[Formalization strategy for Lemma~\ref{lem:low-rank-support}]
This lemma can be formalized either through the apolar Hilbert-function argument above, or as a standalone finite-dimensional theorem about degree-four inverse systems.  The latter statement is sufficient for the main proof:
\[
\rank C_\lambda\le3
\Rightarrow
\exists W,A,
\ E=W\oplus A,\ \dim W\le3,\ AE\subseteq K,
\]
and similarly with \(\dim W\le2\) when \(\rank C_\lambda\le2\), and \(\dim W\le1\) when \(\rank C_\lambda=1\) and \(\lambda\ne0\).
No order or positivity of \(\lambda\) is involved in this structural lemma.
\end{remark}

\subsection{Binary rank at most two: a negative square}

The rank-two and rank-one cases need a negative \emph{pure square} \(x^2\), not merely an arbitrary quadratic with negative square.

\begin{lemma}[Binary low-rank indefinite forms have a negative pure square]\label{lem:binary-negative-square}
Let \(W\) be two-dimensional and let \(\lambda\) be a quartic functional supported on \(W\), so its catalecticant acts on \(\Sym^2W\).  If
\[
\rank C_\lambda\le2
\]
and there exists \(q\in\Sym^2W\) with
\[
\lambda(q^2)<0,
\]
then there exists \(x\in W\) with
\[
\lambda(x^4)<0.
\]
\end{lemma}

\begin{proof}
Choose a basis \(x,y\) of \(W\).  Write
\[
\lambda(x^4)=a,
\quad
\lambda(x^3y)=b,
\quad
\lambda(x^2y^2)=c,
\quad
\lambda(xy^3)=d,
\quad
\lambda(y^4)=e.
\]
In the basis \((x^2,xy,y^2)\), the catalecticant matrix is the Hankel matrix
\[
M=
\begin{pmatrix}
a&b&c\\
b&c&d\\
c&d&e
\end{pmatrix}.
\]
The assumption \(\exists q,\lambda(q^2)<0\) means that \(M\) is not positive semidefinite.

If \(\rank M=1\), the Hankel relations force
\[
M=\rho
\begin{pmatrix}
\alpha^2\\ \alpha\beta\\ \beta^2
\end{pmatrix}
\begin{pmatrix}
\alpha^2&\alpha\beta&\beta^2
\end{pmatrix}
\]
for some real \(\alpha,\beta\).  Since \(M\) is not psd, \(\rho<0\).  Taking \(x'=\alpha x+\beta y\) gives \(\lambda((x')^4)<0\).

If \(\rank M=2\), its kernel is a real one-dimensional space generated by a binary quadratic.  Up to a real linear change of basis, the kernel generator is one of
\[
xy,
\qquad
y^2,
\qquad
x^2+y^2.
\]
The corresponding forms are as follows.

\begin{enumerate}[label=(\roman*)]
\item If \(xy\in\ker M\), then \(b=c=d=0\), so only \(a x^4+e y^4\) remains.  Since \(M\) is not psd, either \(a<0\) or \(e<0\), giving a negative pure fourth power.

\item If \(y^2\in\ker M\), then \(c=d=e=0\), so along \(x+ty\) one has
\[
\lambda((x+ty)^4)=a+4bt.
\]
Rank two forces \(b\ne0\), and some real \(t\) makes this negative.

\item If \(x^2+y^2\in\ker M\), then
\[
a+c=0,
\qquad b+d=0,
\qquad c+e=0.
\]
Thus
\[
\lambda((x+ty)^4)=a(1-6t^2+t^4)+4b(t-t^3).
\]
If \(a\ne0\), then \(t=0\) or \(t=1\) gives opposite signs, since the values are \(a\) and \(-4a\).  If \(a=0\), then \(b\ne0\) by rank two, and \(4b(t-t^3)\) is negative for some real \(t\).
\end{enumerate}
Thus in all cases some linear form has negative fourth-power value.
\end{proof}

\begin{corollary}[Rank-one normal form]\label{cor:rank-one-normal-form}
If \(\rank C_\lambda=1\), \(\lambda\) has a negative square, and \(E\) is four-dimensional, then there is a decomposition
\[
E=\R x\oplus B,
\qquad \dim B=3,
\]
such that
\[
K=xB\oplus\Sym^2B,
\qquad
\lambda(x^4)<0.
\]
\end{corollary}

\begin{proof}
By Lemma~\ref{lem:low-rank-support}, \(\lambda\) is supported on at most one essential variable.  Since \(\lambda\) has a negative square, it is nonzero, so the support is exactly one-dimensional.  Let \(x\) span the support and let \(B\) be a complementary three-dimensional subspace.  Then every quadratic except the \(x^2\)-component lies in the kernel, giving the displayed formula for \(K\).  The negative-square condition gives \(\lambda(x^4)<0\) after scaling \(x\).
\end{proof}

\section{Full-rank proof by catalecticant rank}

Assume now that \(u\) is second-order critical, \(R\ne0\), and the coordinates \(u_i\) are linearly independent.  Then
\[
\dim L=7,
\qquad
L\subseteq K,
\qquad
\rank C_\lambda=k\le3.
\]
By Lemma~\ref{lem:negative-sos-summand}, \(\lambda\) has a negative square.  In particular \(k\ne0\), because if \(k=0\), every quadratic square from the SOS representation of \(p\) would have \(\lambda(q_j^2)=0\), contradicting \(\lambda(p)<0\).

Thus \(k\in\{1,2,3\}\).  Each case produces a syzygy of the form required by Lemma~\ref{lem:syzygy-negative-curvature}.

\subsection{Case \(k=3\)}

Here
\[
\dim K=10-3=7=\dim L,
\]
so
\[
L=K.
\]
By Lemma~\ref{lem:low-rank-support}, choose
\[
E=W\oplus A,
\qquad
\dim A\ge1,
\qquad
AE\subseteq K.
\]
Pick \(0\ne z\in A\).  Then
\[
zE\subseteq K=L,
\]
so in particular
\[
z^2\in L,
\qquad
z\ell\in L\quad(\ell\in W).
\]

Choose a quadratic \(q_0\) with \(\lambda(q_0^2)<0\).  Decompose
\[
q_0=q_W+q_K,
\qquad
q_W\in\Sym^2W,
\qquad
q_K\in AE\subseteq K.
\]
Since \(q_K\in K\), all terms involving \(q_K\) vanish under \(\lambda\), and therefore
\[
\lambda(q_W^2)=\lambda(q_0^2)<0.
\]
Write the quadratic \(q_W\) as a finite sum of products of linear forms in \(W\):
\[
q_W=\sum_{j\in J}\ell_jm_j,
\qquad
\ell_j,m_j\in W.
\]
Now define
\[
M_j=z\ell_j\in L,
\qquad
N_j=zm_j\in K,
\qquad
s=z^2\in L\setminus\{0\},
\qquad
q=q_W.
\]
Then
\[
\sum_{j\in J}M_jN_j
=\sum_j(z\ell_j)(zm_j)
=z^2\sum_j\ell_jm_j
=sq.
\]
Also \(\lambda(q^2)<0\).  Lemma~\ref{lem:syzygy-negative-curvature} gives a negative Hessian direction, contradiction.

\subsection{Case \(k=2\)}

By Lemma~\ref{lem:low-rank-support}, choose a decomposition
\[
E=W\oplus A,
\qquad
\dim W=2,
\qquad
\dim A=2,
\qquad
AE\subseteq K.
\]
The support is genuinely two-dimensional because the catalecticant rank is two.  By Lemma~\ref{lem:binary-negative-square}, after choosing a suitable \(x\in W\), we have
\[
\lambda(x^4)<0.
\]
Set
\[
U=xA\oplus\Sym^2A\subseteq K.
\]
Then
\[
\dim U=2+3=5,
\qquad
\dim K=10-2=8.
\]
Because \(L\subseteq K\), \(\dim L=7\), and \(U\subseteq K\), the Grassmann dimension inequality gives
\[
\dim(L\cap U)
\ge
\dim L+
\dim U-
\dim K
=7+5-8=4.
\]
Put
\[
P=L\cap U.
\]
Inside the direct sum \(U=xA\oplus\Sym^2A\), define
\[
M=\{a\in A: xa\in P\},
\qquad
N=P\cap\Sym^2A.
\]
Then
\[
\dim M=\dim(P\cap xA)
\ge 4+2-5=1,
\]
and
\[
\dim N
\ge 4+3-5=2.
\]
Since \(A\) is two-dimensional and \(M\ne0\), the product space
\[
MA=\Span\{ab:a\in M,b\in A\}\subseteq\Sym^2A
\]
has dimension at least two.  Because \(\dim\Sym^2A=3\),
\[
\dim N+
\dim(MA)
\ge2+2>3,
\]
so
\[
N\cap MA\ne\{0\}.
\]
Choose
\[
0\ne s\in N\cap MA.
\]
Then \(s\in N\subseteq P\subseteq L\), and since \(s\in MA\), we can write
\[
s=\sum_{j\in J}a_jb_j,
\qquad
a_j\in M,
\quad b_j\in A.
\]
For each \(j\), membership \(a_j\in M\) gives
\[
xa_j\in P\subseteq L,
\]
while \(b_j\in A\) gives
\[
xb_j\in xA\subseteq AE\subseteq K.
\]
Now use the identity
\[
\sum_j(xa_j)(xb_j)
=x^2\sum_j a_jb_j
=x^2s.
\]
This is Lemma~\ref{lem:syzygy-negative-curvature} with
\[
M_j=xa_j,
\qquad
N_j=xb_j,
\qquad
q=x^2,
\qquad
s=s.
\]
The negativity condition is \(\lambda(q^2)=\lambda(x^4)<0\).  Hence a negative Hessian direction exists, contradiction.

\subsection{Case \(k=1\)}

By Corollary~\ref{cor:rank-one-normal-form}, choose
\[
E=\R x\oplus B,
\qquad
\dim B=3,
\]
such that
\[
K=xB\oplus\Sym^2B,
\qquad
\lambda(x^4)<0.
\]
The dimensions are
\[
\dim K=9,
\qquad
\dim xB=3,
\qquad
\dim\Sym^2B=6.
\]
Define
\[
M=\{b\in B: xb\in L\},
\qquad
N=L\cap\Sym^2B.
\]
By the dimension inequality inside \(K\),
\[
\dim M
=
\dim(L\cap xB)
\ge 7+3-9=1,
\]
and
\[
\dim N
=
\dim(L\cap\Sym^2B)
\ge 7+6-9=4.
\]
Since \(M\ne0\) and \(B\) is three-dimensional, multiplication by a nonzero element of \(M\) embeds \(B\) into \(\Sym^2B\), so
\[
\dim(MB)
\ge3.
\]
As \(\dim\Sym^2B=6\),
\[
\dim N+
\dim(MB)
\ge4+3>6.
\]
Therefore
\[
N\cap MB\ne\{0\}.
\]
Choose
\[
0\ne s\in N\cap MB.
\]
Then \(s\in L\), and because \(s\in MB\), write
\[
s=\sum_{j\in J}b_jc_j,
\qquad
b_j\in M,
\quad c_j\in B.
\]
Then
\[
xb_j\in L,
\qquad
xc_j\in xB\subseteq K,
\]
and
\[
\sum_j(xb_j)(xc_j)
=x^2\sum_j b_jc_j
=x^2s.
\]
Apply Lemma~\ref{lem:syzygy-negative-curvature} with
\[
q=x^2,
\qquad
\lambda(q^2)=\lambda(x^4)<0.
\]
Again this contradicts second-order criticality.

\section{Main proof}

\begin{theorem}[Rank-seven dehomogenized quaternary quartic theorem]\label{thm:main}
For every positive-definite bilinear form \(B\), every SOS quartic \(p\), and every admissible rank-seven factor \(u\), if \(u\) is second-order critical in the sense of the Lean definition, then
\[
\sigma(u)-p=0.
\]
\end{theorem}

\begin{proof}
Let \(R=\sigma(u)-p\).  Suppose \(R\ne0\).  Define \(\lambda(f)=B(f,R)
\).  Lemma~\ref{lem:lambda-p-negative} gives \(\lambda(p)<0\), and Lemma~\ref{lem:negative-sos-summand} gives a quadratic \(q\) with \(\lambda(q^2)<0\).

Let \(L=\Span\{u_i\}\) and \(K=\{a\in\Quad: \forall b\in\Quad,\lambda(ab)=0\}\).  By Lemma~\ref{lem:L-subset-K}, \(L\subseteq K\).

If the \(u_i\) are linearly dependent, Lemma~\ref{lem:rank-deficient} contradicts second-order criticality.  Therefore they are linearly independent and \(\dim L=7\).  Since \(L\subseteq K\subseteq\Quad\) and \(\dim\Quad=10\), the catalecticant rank satisfies
\[
k=10-
\dim K\le3.
\]
Moreover \(k\ne0\), because \(k=0\) would imply every SOS summand square has \(\lambda(q_j^2)=0\), contradicting \(\lambda(p)<0\).

Thus \(k=1,2,\) or \(3\).  The rank-three, rank-two, and rank-one arguments in the previous section each construct a syzygy satisfying Lemma~\ref{lem:syzygy-negative-curvature}.  Each case therefore produces an admissible direction \(v\) with
\[
H_u(v)<0,
\]
contradicting the second-order criticality hypothesis.

The contradiction shows that \(R\ne0\) is impossible.  Hence
\[
R=\sigma(u)-p=0,
\]
which is exactly \lean{residual p u = 0}.
\end{proof}

\section{Suggested Lean lemma inventory}

The following list is meant as a formalization guide.  Names are suggestions only; the target theorem name and statement should remain unchanged.

\subsection*{Finite-dimensional polynomial spaces}

\begin{enumerate}[label=\textbf{F\arabic*.},leftmargin=2.5em]
\item \lean{linSubmodule}, \lean{quadSubmodule}, \lean{quarticSubmodule}: submodules of polynomials of total degree at most \(1,2,4\).
\item \lean{finrank_lin_eq_four}: \(\dim\Lin=4\).
\item \lean{finrank_quad_eq_ten}: \(\dim\Quad=10\).
\item \lean{mul_lin_lin_mem_quad}: product of two affine-linear polynomials is quadratic.
\item \lean{mul_quad_quad_mem_quartic}: product of two quadratics is quartic.
\item \lean{quad_eq_sym2_lin_span}: products of two elements of \(\Lin\) span \(\Quad\).
\end{enumerate}

\subsection*{Residual functional and FOCP}

\begin{enumerate}[label=\textbf{R\arabic*.},leftmargin=2.5em]
\item \lean{lambda_of_residual}: define \(\lambda(f)=B(f,\lean{residual p u})\).
\item \lean{FOCP_self}: from \lean{IsFOCP B p u}, prove \(B(\sigma u,R)=0\).
\item \lean{lambda_p_negative}: if \(R\ne0\), then \(\lambda(p)<0\).
\item \lean{exists_negative_sos_summand}: from \lean{IsSOSQuartic p} and \(\lambda(p)<0\), produce \(q\in\Quad\) with \(\lambda(q^2)<0\).
\item \lean{catalecticantKernel}: define \(K=\{a\in\Quad: \forall b\in\Quad,\lambda(ab)=0\}\).
\item \lean{span_u_subset_catalecticantKernel}: FOCP implies \(L\subseteq K\).
\end{enumerate}

\subsection*{Negative-curvature constructions}

\begin{enumerate}[label=\textbf{N\arabic*.},leftmargin=2.5em]
\item \lean{rankDeficient_negative_curvature}: Lemma~\ref{lem:rank-deficient}.
\item \lean{syzygy_negative_curvature}: Lemma~\ref{lem:syzygy-negative-curvature}.
\item \lean{fullRank_rank_catalecticant_le_three}: if \(\dim L=7\) and \(L\subseteq K\), then \(\rank C_\lambda\le3\).
\end{enumerate}

\subsection*{Low-rank structure}

\begin{enumerate}[label=\textbf{S\arabic*.},leftmargin=2.5em]
\item \lean{lowRankCatalecticant_supportReduction}: Lemma~\ref{lem:low-rank-support}.
\item \lean{binary_rank_le_two_negative_square}: Lemma~\ref{lem:binary-negative-square}.
\item \lean{rank_one_normal_form}: Corollary~\ref{cor:rank-one-normal-form}.
\item \lean{rank_three_syzygy}: construction in the \(k=3\) case.
\item \lean{rank_two_syzygy}: construction in the \(k=2\) case.
\item \lean{rank_one_syzygy}: construction in the \(k=1\) case.
\end{enumerate}

\subsection*{Main theorem}

\begin{enumerate}[label=\textbf{M\arabic*.},leftmargin=2.5em]
\item \lean{QuaternaryQuarticRankSevenNoSpuriousSOCP}: unfold the definition, introduce \(B,p,u\), assume all hypotheses, set \(R=\lean{residual p u}\), and prove by contradiction using Theorem~\ref{thm:main}.
\end{enumerate}

\section{Lean-specific comments}

\begin{itemize}[leftmargin=2em]
\item The proof never needs to modify \lean{QuaternaryQuartic.lean}.  It should live in a separate file, for example \lean{QuaternaryQuarticProof.lean}, importing the statement file.

\item Since \lean{DotForm} is not assumed symmetric in the target, avoid any use of symmetry of \(B\).  Work with \(\lambda(f)=B(f,R)\) in the first argument only.  The catalecticant \(C_\lambda(a,b)=\lambda(ab)\) is symmetric because polynomial multiplication is commutative, not because \(B\) is symmetric.

\item When proving Hessian negativity, use the Lean definition
\[
\lean{hessianTerm B p u v}=B(\sigma(v),R)+2B(A(u,v),A(u,v)).
\]
All constructed directions satisfy \(A(u,v)=0\), so the second term is exactly zero.

\item The finite sums in \lean{sigma} and \lean{A} are over \lean{Fin 7}.  For rank-deficient directions, encode the dependence relation as a coefficient vector \(c:\lean{Fin 7}\to\R\) with \(\sum_i c_i u_i=0\).  The direction is \(v_i=c_i q\).

\item For syzygy directions, if
\[
m_j=\sum_i\alpha_{j,i}u_i,
\qquad
s=\sum_i\beta_i u_i,
\]
define
\[
v_i=\sum_j\alpha_{j,i}n_j-\beta_iq.
\]
The key simplifications are
\[
A(u,v)=0,
\qquad
\lambda(\sigma(v))=\left(\sum_i\beta_i^2\right)\lambda(q^2).
\]
The second identity follows because every \(n_j\in K\) kills every quadratic factor under \(\lambda\).

\item In the dimension arguments, repeatedly use
\[
\dim(U\cap V)\ge\dim U+
\dim V-
\dim W
\]
for subspaces \(U,V\subseteq W\).  Mathlib has the corresponding formula via \(\dim(U+V)+\dim(U\cap V)=\dim U+
\dim V\).

\item The rank-three case is the cleanest case to formalize after the syzygy lemma: since \(\dim K=7=\dim L\), one has \(L=K\), and the syzygy is simply
\[
\sum_j(z\ell_j)(zm_j)=z^2q.
\]

\item The rank-two and rank-one cases are pure finite-dimensional incidence arguments inside \(\Sym^2A\) or \(\Sym^2B\).  They should be formalized as separate lemmas about subspaces of a two- or three-dimensional vector space before being connected back to polynomials.
\end{itemize}

\end{document}
